El presente trabajo tuvo como alcance la reingeniería mecánica parcial y la automatización de un brazo robótico cartesiano para su integración en una línea de empaquetado del Laboratorio de Diseño de Procesos de la Universidad del Valle de Guatemala. El proyecto se enfocó en mejorar el desempeño mecánico del sistema mediante el rediseño de los ejes X, Y y Z, la incorporación de ejes guía y la integración de rodamientos lineales que permitieron obtener un movimiento más fluido, estable y repetible durante la operación.

Asimismo, se contempló la automatización del brazo para ejecutar tareas de manipulación y transferencia de productos mediante una lógica de tipo \textit{pick and place}. Para ello, el sistema identificó objetos dentro del área de trabajo mediante visión computacional, a partir de características previamente definidas como figura y color. Con base en esta información, se determinó la ubicación correspondiente de cada objeto y se ejecutó el movimiento necesario para trasladarlo hacia su posición asignada dentro de la línea de empaquetado.

El proyecto también incluyó la implementación de un modo de operación manual mediante un control inalámbrico, el cual permitió mover los ejes del brazo de forma independiente durante tareas de posicionamiento, pruebas, mantenimiento o intervención del operador. Además, se integró un sistema de calibración automática apoyado en finales de carrera, con el propósito de establecer referencias iniciales de movimiento y mejorar la seguridad operativa del sistema. De igual forma, se incorporó una pantalla como interfaz local para visualizar información relevante, como el modo de operación, el estado de calibración, la detección de objetos y las acciones ejecutadas por el brazo.

Sin embargo, el proyecto se limitó a la automatización del brazo robótico cartesiano y no contempló el diseño completo de una línea de empaquetado industrial. Las pruebas se realizaron en un entorno académico y controlado, por lo que los resultados obtenidos estuvieron condicionados por las características del laboratorio, los recursos disponibles y las dimensiones del prototipo. En consecuencia, el sistema no fue validado bajo condiciones industriales de producción continua, cargas elevadas o ambientes variables.

El sistema de visión computacional se limitó al reconocimiento de figuras y colores definidos para las pruebas del proyecto. Por lo tanto, no se desarrolló un sistema destinado a identificar cualquier tipo de objeto, material, textura o geometría compleja. Su desempeño estuvo condicionado por factores como la iluminación, la posición de la cámara, el contraste entre los objetos y el fondo, y la estabilidad del entorno visual.

Finalmente, aunque el rediseño mecánico estuvo orientado a mejorar la fluidez y repetibilidad del movimiento, el brazo permaneció condicionado por la estructura base existente, la capacidad de los actuadores, el espacio de trabajo disponible y las tolerancias de fabricación. La validación del sistema se realizó mediante pruebas de precisión, repetibilidad, tiempo de respuesta y funcionamiento general, con el objetivo de comprobar la viabilidad del brazo como una plataforma académica funcional para tareas de clasificación, manipulación y transferencia de productos en una línea de empaquetado controlada.